% --- GLOSSAIRE ET LISTE DES ACRONYMES (sans numérotation) ---
\chapter*{Liste des acronymes}
\addcontentsline{toc}{chapter}{Liste des acronymes}
\label{chap:acronymes}

\begin{description}
    \item[ADR] Architecture Decision Record, fiche de décision architecturale.
    \item[API] Application Programming Interface, interface de programmation.
    \item[C4] Context, Container, Component, Code, modèle de modélisation architecturale à quatre niveaux.
    \item[CDP] Commission de protection des Données Personnelles (Sénégal).
    \item[CEDEAO] Communauté Économique des États de l'Afrique de l'Ouest.
    \item[CI/CD] Continuous Integration / Continuous Deployment, intégration et déploiement continus.
    \item[CMTR] Conseil Malien des Transporteurs Routiers.
    \item[COSEC] Conseil Sénégalais des Chargeurs.
    \item[CRUD] Create, Read, Update, Delete, opérations élémentaires de persistance.
    \item[CSV] Comma-Separated Values, format de fichier tabulaire.
    \item[DAG] Directed Acyclic Graph, graphe orienté acyclique (unité d'orchestration Airflow).
    \item[EF] Exigence Fonctionnelle.
    \item[ENF] Exigence Non Fonctionnelle.
    \item[ERP] Enterprise Resource Planning, progiciel de gestion intégré.
    \item[ETL] Extract, Transform, Load, extraction, transformation, chargement.
    \item[FCFA] Franc de la Communauté Financière Africaine.
    \item[IA] Intelligence Artificielle.
    \item[JWT] JSON Web Token, jeton d'authentification signé.
    \item[LLM] Large Language Model, grand modèle de langage.
    \item[MoSCoW] Must have, Should have, Could have, Won't have, méthode de priorisation des exigences.
    \item[ODD] Objectifs de Développement Durable.
    \item[OPA] Observatoire des Pratiques Anormales (UEMOA).
    \item[PME] Petite et Moyenne Entreprise.
    \item[RAG] Retrieval-Augmented Generation, génération augmentée par récupération.
    \item[RBAC] Role-Based Access Control, contrôle d'accès à base de rôles.
    \item[ROI] Return On Investment, retour sur investissement.
    \item[SPA] Single Page Application, application web à page unique.
    \item[SPL] Software Product Line, ligne de produits logiciels.
    \item[SQL] Structured Query Language, langage d'interrogation de bases de données relationnelles.
    \item[SUS] System Usability Scale, questionnaire standardisé d'utilisabilité perçue.
    \item[TAM] Technology Acceptance Model, modèle d'acceptation de la technologie.
    \item[TLS] Transport Layer Security, protocole de chiffrement des communications.
    \item[UEMOA] Union Économique et Monétaire Ouest-Africaine.
    \item[UML] Unified Modeling Language, langage de modélisation unifié.
    \item[USD] United States Dollar, dollar américain.
    \item[UTAUT] Unified Theory of Acceptance and Use of Technology, théorie unifiée de l'acceptation et de l'utilisation de la technologie.
\end{description}

\chapter*{Glossaire}
\addcontentsline{toc}{chapter}{Glossaire}
\label{chap:glossaire}

\begin{description}
    \item[Corridor Dakar--Bamako] Axe logistique routier d'environ 1~400~km reliant le port de Dakar à Bamako, colonne vertébrale du commerce extérieur malien.
    \item[Embedding] Représentation vectorielle dense d'un texte, permettant d'en comparer le sens par similarité géométrique plutôt que par correspondance de mots-clés.
    \item[Feature model (modèle de caractéristiques)] Représentation arborescente des caractéristiques communes (invariants) et variables (points de variation) d'une ligne de produits logiciels.
    \item[Frugalité (architecture frugale)] Conception logicielle et d'infrastructure minimisant le coût et les ressources consommées, sous une contrainte budgétaire forte et assumée dès la conception.
    \item[Idempotence] Propriété d'une opération qui, répétée sur les mêmes données, produit le même résultat qu'une seule exécution, ici, rejouer le pipeline ETL n'insère jamais de doublon.
    \item[pgvector] Extension de PostgreSQL ajoutant un type de donnée vectoriel et des opérateurs de recherche par similarité, utilisée pour le retrieval du RAG local.
    \item[Point de variation] Endroit de l'architecture où la variabilité entre membres d'une ligne de produits se manifeste et doit être résolu par une configuration ou un choix propre à chaque instance.
    \item[RAG (Retrieval-Augmented Generation)] Technique augmentant un grand modèle de langage par la récupération préalable de passages ou de données pertinentes, avant génération de la réponse.
    \item[Souveraineté des données] Garantie qu'une donnée ne quitte pas une infrastructure maîtrisée ni la juridiction dont elle relève, notamment vis-à-vis d'un service tiers hébergé à l'étranger.
    \item[Text-to-SQL] Technique par laquelle un modèle de langage traduit une question posée en langage naturel en une requête SQL, exécutée ensuite sur la base de données.
\end{description}
